% ============================================================
% §03 — Schema vNext
% Author: John (YAML/Schema Specialist)
% Last updated: 2026-04-18
% ============================================================

\chapter{Schema vNext}

\label{chap:schema}

This chapter is the normative specification for all gert v2 data structures:
runbook documents, tool definitions, and provider definitions.
It supersedes all v1 schema documentation and defines the contracts that Brian's
parser, Ken's runtime, and the VS~Code extension must implement.

The design philosophy is \emph{minimal core, extensible periphery}: the core
language fields that the runtime must understand are small, typed, and stable;
everything else is declared as an extension with an explicit namespace.  Where
v1 accumulated surface area organically, v2 draws a hard boundary and enforces
it at parse time.

% ============================================================
\section{Schema Identity and Versioning}
\label{sec:schema-identity}
% ============================================================

\subsection{apiVersion Values}
\label{subsec:apiversion}

Every gert document carries a mandatory \texttt{apiVersion} field that
identifies both the document kind and the schema version.  The v2 values are:

\begin{center}
\begin{tabular}{lll}
\hline
\textbf{Document kind} & \textbf{v1 value} & \textbf{v2 value} \\
\hline
Runbook                & \texttt{runbook/v1} & \texttt{runbook/v2} \\
Tool definition        & \texttt{tool/v0}    & \texttt{tool/v2}   \\
Provider definition    & \texttt{provider/v0}& \texttt{provider/v2} \\
\hline
\end{tabular}
\end{center}

The format is \texttt{<kind>/<major>}.  Minor and patch changes do not alter
the \texttt{apiVersion} string; they are tracked by the schema artifact version
(see \S\ref{sec:schema-artifacts}).

\subsection{Self-Description via \$schema}
\label{subsec:self-description}

Following JSON~Schema~\cite{jsonschema-spec} best practice and the pattern used by GitHub Actions,
Argo~Workflows, and Kubernetes~CRDs, every v2 document \emph{should} carry a
\texttt{\$schema} field pointing to the canonical JSON~Schema URL:

\begin{verbatim}
$schema: "https://schemas.gert.dev/runbook/v2.json"
apiVersion: runbook/v2
id: service-health
name: Service Health Check
...
\end{verbatim}

The \texttt{\$schema} field is optional at runtime (the runtime uses
\texttt{apiVersion} for dispatch) but is \emph{strongly recommended} because:

\begin{itemize}
  \item VS~Code and JetBrains IDEs provide schema-backed auto-completion and
        inline validation when \texttt{\$schema} is present.
  \item The schema URL carries an implicit version contract: updating the URL
        updates the validation rules visible to tooling.
  \item It enables offline validation with any JSON~Schema Draft~2020-12
        validator (ajv, jsonschema, check-jsonschema).
\end{itemize}

Schema URLs for each document kind:

\begin{verbatim}
https://schemas.gert.dev/runbook/v2.json
https://schemas.gert.dev/tool/v2.json
https://schemas.gert.dev/provider/v2.json
\end{verbatim}

\subsection{Versioning Strategy}
\label{subsec:versioning-strategy}

gert schema versioning follows a simplified semantic versioning model applied
to the \texttt{apiVersion} major component:

\begin{itemize}
  \item \textbf{Major increment} (e.g.\ \texttt{runbook/v1} → \texttt{runbook/v2}):
        A breaking change.  The parser \emph{must not} silently accept the old
        format.  Automated migration is provided via \texttt{gert migrate}.
  \item \textbf{Minor/patch changes within a major}: Additive, backward-compatible.
        New optional fields may be added.  Existing fields may not be removed or
        renamed.  The JSON~Schema artifact version tracks these with a semver
        suffix (e.g.\ \texttt{runbook/v2.json} at artifact version
        \texttt{2.1.0}).
\end{itemize}

\paragraph{What constitutes a breaking change}

A breaking change is any of:
\begin{itemize}
  \item Renaming or removing a required field.
  \item Narrowing an existing enum (removing a valid value).
  \item Changing a field type (e.g.\ string → object).
  \item Changing the meaning of a field such that existing runbooks produce
        different runtime behavior.
  \item Removing a step type.
\end{itemize}

Adding new optional fields, adding new enum values, and loosening existing
constraints are \emph{not} breaking changes.

\subsection{Compatibility Contract: v1 → v2}
\label{subsec:compatibility-contract}

\begin{center}
\begin{tabular}{llp{6cm}}
\hline
\textbf{v1 field} & \textbf{v2 status} & \textbf{Notes} \\
\hline
\texttt{apiVersion: runbook/v0} & Removed & Handled by \texttt{gert migrate} \\
\texttt{apiVersion: runbook/v1} & Deprecated & Compatibility mode in v2 parser \\
\texttt{meta.name}              & Promoted  & Becomes top-level \texttt{name} \\
\texttt{meta.kind}              & Promoted  & Becomes top-level \texttt{kind} \\
\texttt{meta.description}       & Promoted  & Becomes top-level \texttt{description} \\
\texttt{meta.vars}              & Retained  & Moved to top-level \texttt{vars} \\
\texttt{meta.inputs}            & Retained  & Moved to top-level \texttt{inputs} \\
\texttt{meta.governance}        & Retained  & Moved to top-level \texttt{governance} \\
\texttt{meta.prose}             & Retained  & Moved to top-level \texttt{prose} \\
\texttt{meta.defaults}          & Retained  & Moved to top-level \texttt{defaults} \\
\texttt{imports}                & Retained  & Unchanged \\
\texttt{tools}                  & Replaced  & Replaced by \texttt{toolRefs} (typed) \\
\texttt{flow}                   & Retained  & Canonical execution root \\
\texttt{tree}                   & Removed   & Use \texttt{flow} \\
\texttt{steps}                  & Removed   & Flat step list was never canonical \\
\texttt{step.type: collector}   & Renamed   & Becomes \texttt{type: manual} \\
\texttt{step.type: router}      & Renamed   & Becomes \texttt{type: end} \\
\texttt{step.type: choice}      & Retained  & Unchanged \\
\texttt{step.type: noop}        & Removed   & Use \texttt{when:} guard on any step \\
\texttt{step.type: extension}   & Retained  & Unchanged \\
\texttt{step.with.argv}         & Deprecated& Use \texttt{type: cli} with \texttt{command:} \\
\texttt{step.assertions}        & Deprecated& Use \texttt{step.type: assert} \\
\hline
\end{tabular}
\end{center}

The v2 parser reads \texttt{runbook/v1} documents in \textbf{compatibility mode}:
all v1 field paths are accepted and silently normalized to their v2 equivalents
before validation.  Running \texttt{gert migrate} permanently rewrites the file
to v2 canonical form.  v2 will not support \texttt{runbook/v0}; documents at
that version must be migrated to v1 first or directly to v2 via a two-stage
migration.

\subsection{\texttt{gert migrate} Behavior}
\label{subsec:migrate}

\begin{verbatim}
# Dry-run preview
gert migrate --dry-run runbook.yaml

# In-place rewrite
gert migrate runbook.yaml

# Migrate a whole directory
gert migrate ./runbooks/
\end{verbatim}

Migration rules applied automatically:
\begin{enumerate}
  \item Rewrite \texttt{apiVersion} to \texttt{runbook/v2}.
  \item Add \texttt{\$schema} field.
  \item Promote \texttt{meta.name}, \texttt{meta.kind}, \texttt{meta.description}
        to top-level.
  \item Move remaining \texttt{meta.*} fields to their v2 top-level equivalents.
  \item Rename \texttt{tree:} to \texttt{flow:}.
  \item Rename \texttt{step.type: collector} → \texttt{manual},
        \texttt{step.type: router} → \texttt{end}.
  \item Rewrite \texttt{tools: [name]} shorthand to \texttt{toolRefs:} typed
        declarations.
  \item Add \texttt{id:} to runbook if absent (derived from \texttt{name}).
\end{enumerate}

% ============================================================
\section{Core Runbook Fields (v2)}
\label{sec:core-fields}
% ============================================================

\subsection{Top-Level Field Inventory}
\label{subsec:field-inventory}

The v2 runbook document has the following top-level fields:

\begin{center}
\begin{tabular}{lllp{5.5cm}}
\hline
\textbf{Field} & \textbf{Type} & \textbf{Required} & \textbf{Description} \\
\hline
\texttt{\$schema}     & string  & No  & JSON~Schema URL (recommended) \\
\texttt{apiVersion}   & string  & Yes & Must be \texttt{runbook/v2} \\
\texttt{id}           & string  & Yes & Unique machine identifier (slug) \\
\texttt{name}         & string  & Yes & Human-readable display name \\
\texttt{kind}         & enum    & No  & \texttt{mitigation | reference | composable | rca} \\
\texttt{description}  & string  & No  & Short prose summary \\
\texttt{vars}         & map     & No  & Default variable values \\
\texttt{inputs}       & map     & No  & Input declarations (\S\ref{sec:inputs-outputs}) \\
\texttt{outputs}      & map     & No  & Output declarations (\S\ref{sec:inputs-outputs}) \\
\texttt{toolRefs}     & array   & No  & Tool references (\S\ref{subsec:toolrefs}) \\
\texttt{imports}      & map     & No  & Named runbook imports \\
\texttt{defaults}     & object  & No  & Default step settings \\
\texttt{governance}   & object  & No  & Governance policy \\
\texttt{prose}        & object  & No  & Human-readable TSG sections \\
\texttt{metadata}     & map     & No  & Arbitrary key-value annotations \\
\texttt{flow}         & array   & Yes & Execution tree (TreeNode array) \\
\hline
\end{tabular}
\end{center}

\paragraph{id field}
The \texttt{id} is a slug: lowercase alphanumeric with hyphens, matching
\texttt{\^{}[a-z][a-z0-9-]*[a-z0-9]\$}.  It uniquely identifies a runbook within
a project and is used by the \texttt{invoke} step and the trace record.  It must
be stable across renames; changing it is a breaking change for any runbook that
invokes it by id.

\paragraph{kind values}
\begin{itemize}
  \item \texttt{mitigation} — Incident response procedure; governance defaults
        to \texttt{require-approval} for destructive effects.
  \item \texttt{reference} — Documentation-first runbook; dry-run by default.
  \item \texttt{composable} — Sub-runbook intended for invoke from other
        runbooks; not directly executable from CLI.
  \item \texttt{rca} — Root-cause analysis template.
\end{itemize}

\subsection{toolRefs}
\label{subsec:toolrefs}

In v1, tools were declared as a flat string array (\texttt{tools: [curl, nslookup]}).
In v2, \texttt{toolRefs} is a typed array that allows aliasing and explicit path
overrides:

\begin{verbatim}
toolRefs:
  - name: curl
  - name: nslookup
  - name: k8s
    path: ./tools/k8s.tool.yaml   # override default discovery
    alias: kubectl-wrapper        # optional local alias
\end{verbatim}

The default discovery path for \texttt{name: foo} is \texttt{tools/foo.tool.yaml}
relative to the runbook file.  The \texttt{path} field overrides this.
The \texttt{alias} field lets the runbook reference the tool by a local name
that differs from the tool's declared \texttt{meta.name}.

\subsection{Minimal Valid v2 Runbook}
\label{subsec:minimal-runbook}

\begin{verbatim}
$schema: "https://schemas.gert.dev/runbook/v2.json"
apiVersion: runbook/v2
id: hello-world
name: Hello World

flow:
  - step:
      id: greet
      type: cli
      title: Print greeting
      command: echo
      args: ["Hello, world!"]
      capture:
        output: stdout
  - step:
      id: done
      type: end
      title: Complete
      outcome:
        category: resolved
        code: success
\end{verbatim}

\subsection{Complete Annotated Example}
\label{subsec:full-example}

\begin{verbatim}
$schema: "https://schemas.gert.dev/runbook/v2.json"
apiVersion: runbook/v2
id: service-health-check
name: Service Health Check
kind: mitigation
description: |
  Validate DNS resolution and HTTP endpoint reachability
  for a named service. Escalates on failure.

vars:
  default_timeout: "30s"

inputs:
  hostname:
    type: string
    required: true
    description: "Hostname to check"
    from: prompt
  incident_id:
    type: string
    required: false
    description: "Incident tracking ID"
    from: incident.id       # provider-resolved

outputs:
  health_status:
    type: string
    description: "Final health assessment"
    from: captures.http_result

toolRefs:
  - name: curl
  - name: nslookup

governance:
  rules:
    - effects: ["network.probe"]
      action: allow
    - effects: ["service.restart"]
      action: require-approval
      min_approvers: 2
  redact:
    - pattern: "(?i)Authorization:\\s*\\S+"
      replace: "Authorization: <redacted>"

defaults:
  timeout: "60s"

prose:
  safety: |
    Run this runbook only during an active incident. Do not
    loop — each execution creates a trace record.
  prerequisites: "Requires curl >= 7.80 and nslookup on PATH."

metadata:
  team: platform-sre
  runbook-url: "https://wiki.example.com/runbooks/service-health"
  severity: P2

flow:
  - step:
      id: check_dns
      type: tool
      title: "Resolve DNS: {{ .hostname }}"
      tool:
        name: nslookup
        action: lookup
        args:
          host: "{{ .hostname }}"
      capture:
        dns_result: stdout
      contract:
        effects: ["network.probe"]
        reads: ["hostname"]
        writes: ["dns_result"]

  - step:
      id: check_http
      type: tool
      title: "HTTP health check: {{ .hostname }}"
      tool:
        name: curl
        action: head
        args:
          url: "https://{{ .hostname }}/healthz"
      capture:
        http_result: stdout
      timeout: "{{ .default_timeout }}"

  - step:
      id: evaluate
      type: branch
      title: Evaluate health results
      branches:
        - condition: '{{ contains .http_result "200" }}'
          label: Healthy
          steps:
            - step:
                id: resolved
                type: end
                title: Service is healthy
                outcome:
                  category: resolved
                  code: http_ok
        - condition: '{{ not (contains .http_result "200") }}'
          label: Unhealthy
          steps:
            - step:
                id: escalate
                type: end
                title: Escalate to on-call
                outcome:
                  category: escalated
                  code: http_fail
\end{verbatim}

% ============================================================
\section{Input and Output Declarations}
\label{sec:inputs-outputs}
% ============================================================

\subsection{Input Declarations}
\label{subsec:inputs}

Inputs declare variables that must be resolved before execution begins.  In v2
the \texttt{inputs} map is a top-level field (promoted from \texttt{meta.inputs}).

\begin{verbatim}
inputs:
  <name>:
    type: string | number | boolean | secret
    required: true | false          # default: true
    description: "..."
    default: "..."                  # valid only when required: false
    example: "..."
    pattern: "^[a-z]+"             # regex validation (string only)
    from: <binding>                 # resolution source
\end{verbatim}

\paragraph{type field}
The \texttt{type} field (new in v2) declares the value's type.  Supported values:

\begin{itemize}
  \item \texttt{string} — Default. UTF-8 text.
  \item \texttt{number} — Integer or float, passed as string in templates.
  \item \texttt{boolean} — \texttt{true} or \texttt{false}.
  \item \texttt{secret} — String, redacted in traces and log output.
\end{itemize}

\paragraph{from: binding syntax}

The \texttt{from} field specifies how the value is resolved.  v2 binding
forms:

\begin{center}
\begin{tabular}{lp{7.5cm}}
\hline
\textbf{Binding} & \textbf{Behavior} \\
\hline
\texttt{prompt}          & Ask the operator interactively at runtime \\
\texttt{env.<NAME>}      & Read from environment variable \texttt{NAME} \\
\texttt{file.<path>}     & Read first line from a local file \\
\texttt{<provider>.<field>} & Delegate to a registered input provider (e.g.\ \texttt{incident.id}) \\
\texttt{var.<name>}      & Copy from another variable or capture \\
\hline
\end{tabular}
\end{center}

If \texttt{from} is omitted, the field must have a \texttt{default} value or the
runtime will reject the runbook at semantic validation time.

\subsection{Output Declarations}
\label{subsec:outputs}

Outputs declare values that a composable runbook exposes to its caller via the
\texttt{invoke} step's output extraction.

\begin{verbatim}
outputs:
  <name>:
    type: string | number | boolean | secret
    description: "..."
    from: captures.<captureName>    # must reference a capture variable
\end{verbatim}

The \texttt{from: captures.*} binding is the only supported form in v2.
The named capture must be written by at least one step in the runbook; the
semantic validator checks this.

\subsection{Expression Language}
\label{subsec:expressions}

v2 retains \textbf{Go templates} (\texttt{text/template}) as the expression
language for all interpolated fields (\texttt{title}, \texttt{args}, \texttt{when},
\texttt{condition}, \texttt{instructions}, etc.).  The decision to retain Go
templates is grounded in two facts:
\begin{enumerate}
  \item The v1 template engine is battle-tested and broadly used in v1 runbooks.
  \item Go templates' pipelines and custom function maps provide sufficient
        expressiveness for operational conditionals.
\end{enumerate}

The template data context is the \emph{run variable map}: all \texttt{vars},
resolved \texttt{inputs}, and accumulated \texttt{captures} merged into a single
flat \texttt{map[string]string}.  Access by dot notation:
\texttt{\{\{\ .hostname\ \}\}}.

\paragraph{Built-in template functions (v2 additions)}

v2 adds a standard function library accessible in all templates:

\begin{center}
\begin{tabular}{ll}
\hline
\textbf{Function} & \textbf{Signature} \\
\hline
\texttt{contains}     & \texttt{(s, substr string) bool} \\
\texttt{hasPrefix}    & \texttt{(s, prefix string) bool} \\
\texttt{hasSuffix}    & \texttt{(s, suffix string) bool} \\
\texttt{trim}         & \texttt{(s string) string} \\
\texttt{toLower}      & \texttt{(s string) string} \\
\texttt{toUpper}      & \texttt{(s string) string} \\
\texttt{split}        & \texttt{(s, sep string) []string} \\
\texttt{join}         & \texttt{(elems []string, sep string) string} \\
\texttt{env}          & \texttt{(name string) string} \\
\texttt{now}          & \texttt{() string} (RFC3339) \\
\texttt{runID}        & \texttt{() string} \\
\texttt{not}          & \texttt{(v bool) bool} \\
\texttt{eq, ne, lt, gt, le, ge} & Standard comparison \\
\hline
\end{tabular}
\end{center}

All functions from \texttt{text/template} (built-ins: \texttt{and}, \texttt{or},
\texttt{index}, \texttt{len}, \texttt{printf}, etc.) remain available.

\paragraph{Template errors}

In v2, a template evaluation error in a \texttt{condition} or \texttt{when}
field is a hard step failure (not a soft skip).  The step transitions to
\texttt{failed} and execution halts unless \texttt{continue\_on\_fail: true} is
set on the step.

% ============================================================
\section{Step Types (v2 Inventory)}
\label{sec:step-types}
% ============================================================

Each step in \texttt{flow} is a \texttt{TreeNode}: either a \texttt{step} object
(with optional \texttt{branches} and \texttt{iterate} siblings) or a standalone
\texttt{iterate} block.

\subsection{Common Step Fields}
\label{subsec:common-fields}

The following fields are available on every step type:

\begin{center}
\begin{tabular}{lllp{5.5cm}}
\hline
\textbf{Field} & \textbf{Type} & \textbf{Required} & \textbf{Description} \\
\hline
\texttt{id}              & string  & Yes & Unique within the runbook \\
\texttt{type}            & enum    & Yes & Step type (see below) \\
\texttt{title}           & string  & No  & Template; displayed in TUI/trace \\
\texttt{subtitle}        & string  & No  & Secondary template annotation \\
\texttt{when}            & string  & No  & Guard expression; skips step if false \\
\texttt{timeout}         & string  & No  & Duration (e.g.\ \texttt{30s}); overrides default \\
\texttt{delay}           & string  & No  & Pre-execution pause (e.g.\ \texttt{5s}) \\
\texttt{retry}           & object  & No  & Retry config (\S\ref{subsec:retry}) \\
\texttt{continue\_on\_fail} & bool & No  & Continue execution on failure \\
\texttt{scope}           & string  & No  & Named scope for variable isolation \\
\texttt{export}          & string array & No & Variables to export from scope \\
\texttt{capture}         & map     & No  & Output capture bindings \\
\texttt{contract}        & object  & No  & Behavioral contract for governance \\
\hline
\end{tabular}
\end{center}

\paragraph{Retry configuration}
\label{subsec:retry}

\begin{verbatim}
retry:
  max: 3              # maximum retry attempts (required)
  interval: "5s"      # initial wait between attempts
  backoff: linear     # linear | exponential (default: linear)
\end{verbatim}

\paragraph{Step contract}

The \texttt{contract} field annotates a step with its behavioral properties for
governance policy evaluation:

\begin{verbatim}
contract:
  effects:
    - network.probe
    - service.restart
  reads:   [hostname, port]
  writes:  [http_result]
  deterministic: true
  idempotent: false
  secrets: [auth_token]
\end{verbatim}

Effects are free-form strings matched against \texttt{governance.rules[].effects}.
The \texttt{secrets} list names variables whose values must be redacted in traces.

\subsection{Step Type: cli}
\label{subsec:step-cli}

Executes a local command directly.  In the v2 runtime \texttt{cli} is the
canonical command step; the v1 \texttt{with.argv} form is accepted in
compatibility mode and rewritten to \texttt{cli} by the migration tool.

\begin{center}
\begin{tabular}{lllp{5cm}}
\hline
\textbf{Field} & \textbf{Type} & \textbf{Required} & \textbf{Description} \\
\hline
\texttt{command}  & string       & Yes & Executable name or full path \\
\texttt{args}     & string array & No  & Arguments (templates expanded) \\
\texttt{env}      & map          & No  & Additional environment variables \\
\texttt{workdir}  & string       & No  & Working directory \\
\texttt{stdin}    & string       & No  & Template string piped to stdin \\
\texttt{capture}  & map          & No  & \texttt{stdout | stderr | exit\_code} → var \\
\hline
\end{tabular}
\end{center}

Governance checks against the \texttt{allowed\_commands} / \texttt{denied\_commands}
lists and any matching \texttt{governance.rules[].effects} are evaluated
before the command executes.  If the rule action is \texttt{require-approval},
the step pauses for human approval before proceeding.

\begin{verbatim}
- step:
    id: tail_logs
    type: cli
    title: "Tail last 100 lines from {{ .log_file }}"
    command: tail
    args: ["-n", "100", "{{ .log_file }}"]
    capture:
      output: stdout
      exit: exit_code
    contract:
      effects: [filesystem.read]
      reads: [log_file]
      writes: [output]
\end{verbatim}

\subsection{Step Type: manual}
\label{subsec:step-manual}

Prompts a human operator for evidence, approvals, or free-form input.  Renamed
from \texttt{collector} in v1.  A \texttt{manual} step does not execute any
command; it pauses execution and waits for the operator to provide the declared
evidence items.

\begin{center}
\begin{tabular}{lllp{5cm}}
\hline
\textbf{Field} & \textbf{Type} & \textbf{Required} & \textbf{Description} \\
\hline
\texttt{instructions}       & string  & No  & Markdown prompt displayed to operator \\
\texttt{required\_evidence} & array   & No  & Evidence items (\S\ref{subsec:evidence}) \\
\texttt{approvals}          & object  & No  & Approval gate config \\
\texttt{choices}            & object  & No  & Radio-button choice selection \\
\hline
\end{tabular}
\end{center}

\paragraph{Evidence items}
\label{subsec:evidence}

\begin{verbatim}
required_evidence:
  - kind: text
    name: summary_notes
  - kind: checklist
    name: preflight
    items:
      - "Verified monitoring dashboard"
      - "Confirmed no other incidents in flight"
  - kind: attachment
    name: screenshot
\end{verbatim}

\paragraph{Approval gate}

\begin{verbatim}
approvals:
  min: 2
  roles: [DRI, change-manager]
  timeout: 4h
  on_timeout: escalate          # escalate | fail | skip
  escalate_to: [manager]
\end{verbatim}

The \texttt{timeout} and \texttt{on\_timeout} fields are new in v2, addressing
the SLA enforcement gap identified in the research brief.

\begin{verbatim}
- step:
    id: validate_dashboard
    type: manual
    title: Validate monitoring dashboard
    instructions: |
      Check the Grafana dashboard for **{{ .hostname }}**.
      Confirm error rate is below 1% and P99 < 500 ms.
    required_evidence:
      - kind: checklist
        name: dashboard_check
        items:
          - "Error rate < 1%"
          - "P99 latency < 500ms"
    approvals:
      min: 1
      roles: [DRI]
      timeout: 30m
      on_timeout: fail
\end{verbatim}

\subsection{Step Type: tool}
\label{subsec:step-tool}

Invokes a named tool action.  The tool must be declared in \texttt{toolRefs}.
The runtime dispatches to the tool's configured transport (stdio, jsonrpc, or
mcp).

\begin{center}
\begin{tabular}{lllp{5cm}}
\hline
\textbf{Field} & \textbf{Type} & \textbf{Required} & \textbf{Description} \\
\hline
\texttt{tool.name}    & string & Yes & Tool name (must match toolRefs entry) \\
\texttt{tool.action}  & string & Yes & Action declared in tool definition \\
\texttt{tool.args}    & map    & No  & Argument key-value map (templates expanded) \\
\texttt{tool.version} & string & No  & Optional version pin \\
\texttt{capture}      & map    & No  & \texttt{stdout | stderr | json.<path>} → var \\
\hline
\end{tabular}
\end{center}

Output capture in v2 adds a \texttt{json.<path>} form for JSON-structured tool
outputs, eliminating the need for follow-up \texttt{cli} steps with \texttt{jq}:

\begin{verbatim}
capture:
  pod_name: json.items[0].metadata.name
  status:   json.items[0].status.phase
\end{verbatim}

\begin{verbatim}
- step:
    id: list_pods
    type: tool
    title: "List pods in {{ .namespace }}"
    tool:
      name: kubectl
      action: get-pods
      args:
        namespace: "{{ .namespace }}"
        output: json
    capture:
      pod_count: json.items | length
      first_pod: json.items[0].metadata.name
    contract:
      effects: [k8s.read]
      reads: [namespace]
      writes: [pod_count, first_pod]
\end{verbatim}

\subsection{Step Type: invoke}
\label{subsec:step-invoke}

Invokes another runbook inline as a sub-procedure.  The invoked runbook's
outputs are mapped into the caller's variable space.

\begin{center}
\begin{tabular}{lllp{5cm}}
\hline
\textbf{Field} & \textbf{Type} & \textbf{Required} & \textbf{Description} \\
\hline
\texttt{invoke.runbook}  & string & Yes & Import alias or relative path \\
\texttt{invoke.inputs}   & map    & No  & Input bindings (template-expanded values) \\
\texttt{invoke.outputs}  & map    & No  & Output extraction: local-var → output-name \\
\texttt{invoke.gate}     & object & No  & Stop-if conditions on child outcome \\
\hline
\end{tabular}
\end{center}

\begin{verbatim}
imports:
  connectivity: ./connectivity-check.runbook.yaml

flow:
  - step:
      id: run_connectivity
      type: invoke
      title: Run connectivity check
      invoke:
        runbook: connectivity
        inputs:
          service_name: "{{ .hostname }}"
          incident_id:  "{{ .incident_id }}"
        outputs:
          conn_status: health_status   # local var ← child output
        gate:
          stop_if: [escalated]
\end{verbatim}

\subsection{Step Type: branch}
\label{subsec:step-branch}

Evaluates a list of conditions and executes the first matching arm.  In v2,
\texttt{branch} is an explicit first-class step type (as well as the inline
\texttt{branches:} sibling on any step node).

\begin{center}
\begin{tabular}{lllp{5cm}}
\hline
\textbf{Field} & \textbf{Type} & \textbf{Required} & \textbf{Description} \\
\hline
\texttt{branches}           & array  & Yes & List of Branch objects \\
\texttt{branches[].condition} & string & Yes & Go template evaluating to \texttt{true}/\texttt{false} \\
\texttt{branches[].label}   & string & No  & Display label for diagram/trace \\
\texttt{branches[].steps}   & array  & No  & Steps to execute if condition matches \\
\hline
\end{tabular}
\end{center}

At most one branch arm executes.  If no arm matches and no arm has an
\texttt{else: true} marker, execution continues past the branch node with no
steps run.  The \texttt{else: true} marker (new in v2) designates a catch-all
arm:

\begin{verbatim}
- step:
    id: route_on_env
    type: branch
    title: Route by environment
    branches:
      - condition: '{{ eq .environment "production" }}'
        label: Production path
        steps:
          - step:
              id: require_approval
              type: manual
              title: Require change-manager approval
              approvals:
                min: 1
                roles: [change-manager]
      - condition: '{{ eq .environment "staging" }}'
        label: Staging path
        steps:
          - step:
              id: log_staging
              type: cli
              title: Log staging deploy
              command: echo
              args: ["Deploying to staging"]
      - else: true
        label: Unknown environment — abort
        steps:
          - step:
              id: abort
              type: end
              title: Unknown environment
              outcome:
                category: escalated
                code: unknown_env
\end{verbatim}

\subsection{Step Type: iterate}
\label{subsec:step-iterate}

Repeats a block of steps over a collection or until a convergence condition
is met.  In v2 \texttt{iterate} is a TreeNode sibling (not a step type in the
\texttt{step.type} enum) — it appears at the tree level alongside \texttt{step}
nodes.

\begin{center}
\begin{tabular}{lllp{5cm}}
\hline
\textbf{Field} & \textbf{Type} & \textbf{Required} & \textbf{Description} \\
\hline
\texttt{iterate.over}  & string & * & Template resolving to comma-separated list \\
\texttt{iterate.as}    & string & No & Loop variable name (default: \texttt{item}) \\
\texttt{iterate.max}   & int    & * & Max iterations (convergence mode) \\
\texttt{iterate.until} & string & * & Convergence condition expression \\
\texttt{iterate.collect} & map  & No & Accumulate captures across passes \\
\texttt{iterate.steps} & array  & Yes & Steps to execute each pass \\
\hline
\multicolumn{4}{l}{* \texttt{over} or (\texttt{max} + \texttt{until}) required; mutually exclusive.} \\
\end{tabular}
\end{center}

The loop variable \texttt{\{\{\ .iteration\ \}\}} (0-indexed) is always available.
In list mode, \texttt{\{\{\ .item\ \}\}} (or the \texttt{as} name) holds the
current element.

\begin{verbatim}
# List iteration
- iterate:
    id: check_all_services
    over: "{{ .service_list }}"   # e.g. "api,worker,scheduler"
    as: svc
    collect:
      results: "{{ .svc }}:{{ .http_result }}"
    steps:
      - step:
          id: probe
          type: tool
          title: "Check {{ .svc }}"
          tool:
            name: curl
            action: head
            args:
              url: "https://{{ .svc }}.internal/healthz"
          capture:
            http_result: stdout

# Convergence iteration
- iterate:
    id: wait_for_deploy
    max: 10
    until: '{{ eq .deploy_status "Running" }}'
    steps:
      - step:
          id: poll_deploy
          type: tool
          title: "Poll deployment (pass {{ .iteration }})"
          tool:
            name: kubectl
            action: get-deploy-status
            args:
              name: "{{ .deployment }}"
          capture:
            deploy_status: json.status.phase
      - step:
          id: wait
          type: cli
          title: Wait 15s
          command: sleep
          args: ["15"]
\end{verbatim}

\subsection{Step Type: parallel}
\label{subsec:step-parallel}

Executes a set of independent step groups concurrently, then joins.  New in v2,
addressing the fan-out/fan-in gap identified in Dennis's research brief.

\begin{center}
\begin{tabular}{lllp{5cm}}
\hline
\textbf{Field} & \textbf{Type} & \textbf{Required} & \textbf{Description} \\
\hline
\texttt{branches}          & array  & Yes & Each branch is an independent group \\
\texttt{join.wait\_for}    & enum   & No  & \texttt{all | any | majority} (default: \texttt{all}) \\
\texttt{join.on\_failure}  & enum   & No  & \texttt{fail | continue} (default: \texttt{fail}) \\
\hline
\end{tabular}
\end{center}

Parallel branches are isolated from each other's variable writes during
execution.  After the join, all \texttt{capture} values from all branches are
merged into the caller's variable map.  Write conflicts (two branches writing the
same variable) are resolved by last-writer-wins and a warning is emitted.

\begin{verbatim}
- step:
    id: parallel_checks
    type: parallel
    title: Run DNS and HTTP checks concurrently
    branches:
      - label: DNS check
        steps:
          - step:
              id: dns_check
              type: tool
              tool:
                name: nslookup
                action: lookup
                args:
                  host: "{{ .hostname }}"
              capture:
                dns_result: stdout
      - label: HTTP check
        steps:
          - step:
              id: http_check
              type: tool
              tool:
                name: curl
                action: head
                args:
                  url: "https://{{ .hostname }}/healthz"
              capture:
                http_result: stdout
    join:
      wait_for: all
      on_failure: fail
\end{verbatim}

\subsection{Step Type: assert}
\label{subsec:step-assert}

Evaluates a list of assertions against captured variables.  Fails the run if
any assertion does not pass.

\begin{verbatim}
- step:
    id: verify_deployment
    type: assert
    title: Verify deployment succeeded
    assert:
      - type: contains
        subject: "{{ .deploy_output }}"
        expected: "successfully rolled out"
      - type: eq
        subject: "{{ .exit_code }}"
        expected: "0"
      - type: json_path
        subject: "{{ .api_response }}"
        path: "$.status"
        expected: "active"
\end{verbatim}

Assertion types: \texttt{contains}, \texttt{not\_contains}, \texttt{matches}
(regex), \texttt{eq}, \texttt{ne}, \texttt{lt}, \texttt{gt}, \texttt{exit\_code},
\texttt{json\_path}.

\subsection{Step Type: compensate}
\label{subsec:step-compensate}

New in v2.  Registers a compensation (rollback) action that is executed in
reverse order if the run fails or is cancelled after this step.  Implements the
saga pattern for long-running operational transactions.

\begin{center}
\begin{tabular}{lllp{5cm}}
\hline
\textbf{Field} & \textbf{Type} & \textbf{Required} & \textbf{Description} \\
\hline
\texttt{compensate.on}   & enum  & No  & \texttt{failure | cancel | any} (default: \texttt{failure}) \\
\texttt{compensate.steps}& array & Yes & Steps to execute as compensation \\
\hline
\end{tabular}
\end{center}

The compensation steps execute in LIFO order relative to the sequence of
\texttt{compensate} nodes encountered during the forward execution pass.
Compensation steps share the same variable context as the point at which the
\texttt{compensate} node was registered.

\begin{verbatim}
- step:
    id: scale_down
    type: tool
    title: Scale down service
    tool:
      name: kubectl
      action: scale
      args:
        deployment: "{{ .deployment }}"
        replicas: "0"

- step:
    id: register_scale_up
    type: compensate
    title: Compensation — scale back up on failure
    compensate:
      on: failure
      steps:
        - step:
            id: scale_up_comp
            type: tool
            title: Restore service scale
            tool:
              name: kubectl
              action: scale
              args:
                deployment: "{{ .deployment }}"
                replicas: "{{ .original_replicas }}"
\end{verbatim}

\subsection{Step Type: end}
\label{subsec:step-end}

Declares a terminal outcome and halts execution.  Renamed from \texttt{router}
in v1.

\begin{verbatim}
- step:
    id: resolved
    type: end
    title: Incident resolved
    outcome:
      category: resolved       # resolved | escalated | no_action | needs_rca
      code: dns_healthy
      meta:
        resolution_time: "{{ .elapsed }}"
\end{verbatim}

% ============================================================
\section{Tool Definition Schema (v2)}
\label{sec:tool-schema}
% ============================================================

Tool definitions describe callable tools: their actions, argument schemas,
transport, and behavioral contract.  In v2 the \texttt{tool/v0} version is
replaced by \texttt{tool/v2}.

\subsection{Tool Definition Field Inventory}
\label{subsec:tool-fields}

\begin{verbatim}
$schema: "https://schemas.gert.dev/tool/v2.json"
apiVersion: tool/v2

meta:
  name: kubectl
  version: "1.2"
  description: Kubernetes CLI wrapper
  binary: kubectl

transport:
  mode: stdio         # stdio | jsonrpc | mcp
  # jsonrpc options:
  # start: ["node", "server.js"]
  # port: 3000
  # mcp options:
  # start: ["python", "mcp_server.py"]

actions:
  get-pods:
    description: List pods in a namespace
    argv: ["get", "pods", "-n", "{{ .namespace }}", "-o", "json"]
    args:
      namespace:
        type: string
        required: true
        description: Kubernetes namespace
      output:
        type: string
        required: false
        default: json
    capture:
      stdout:
        format: json    # json | text | lines
    contract:
      effects: [k8s.read]
      reads: [namespace]
    governance:
      action: allow
\end{verbatim}

\paragraph{New in v2}

\begin{itemize}
  \item \texttt{capture.stdout.format} — declares the output format; enables
        structured \texttt{json.<path>} capture in runbook steps.
  \item \texttt{actions.<name>.contract} — per-action behavioral contract
        automatically merged with any step-level contract override.
  \item \texttt{actions.<name>.governance} — per-action governance policy;
        overrides the runbook-level rules for this action only.
  \item \texttt{meta.version} — explicit tool version for auditing and pinning.
\end{itemize}

% ============================================================
\section{Provider Definition Schema (v2)}
\label{sec:provider-schema}
% ============================================================

Provider definitions declare external input resolution services.
In v2 the \texttt{provider/v0} version is replaced by \texttt{provider/v2}.

\begin{verbatim}
$schema: "https://schemas.gert.dev/provider/v2.json"
apiVersion: provider/v2

meta:
  name: incident
  version: "1.0"
  description: PagerDuty incident context provider
  kind: input-provider    # input-provider | enrichment-provider

transport:
  mode: jsonrpc
  start: ["./providers/pagerduty-provider"]

fields:
  id:
    type: string
    description: Active incident ID
  title:
    type: string
    description: Incident title
  severity:
    type: string
    enum: [P1, P2, P3, P4]
  affected_service:
    type: string
    description: Primary impacted service name

capabilities:
  - input-resolution
\end{verbatim}

The \texttt{fields} map defines the provider's schema.  Each field name
corresponds to a \texttt{from: <provider>.<field>} binding in a runbook input.
The semantic validator checks that all \texttt{from: <provider>.*} bindings
reference a field declared in a loaded provider definition.

% ============================================================
\section{Namespace Convention for Extensions}
\label{sec:namespaces}
% ============================================================

\subsection{The Extension Field Problem}

Strict JSON Schema validation (\texttt{additionalProperties: false}) is a core
v2 design constraint: unknown fields in a runbook must be rejected, not silently
ignored.  At the same time, extensions must be able to attach metadata to steps,
runbooks, and tool definitions without forking the core schema.

\subsection{Convention: \texttt{x-<namespace>:} at the Field Level}

v2 adopts the \textbf{\texttt{x-<namespace>}} prefix convention used by OpenAPI
and GitHub Actions.  Extension fields must:

\begin{enumerate}
  \item Start with \texttt{x-}.
  \item Be followed by a namespace component (lowercase alphanumeric and hyphens,
        matching \texttt{[a-z][a-z0-9-]+}).
  \item Be declared in an \textbf{extension manifest} registered via
        \texttt{gert extension register} or in \texttt{gert.yaml}.
\end{enumerate}

\paragraph{At the runbook level}

\begin{verbatim}
$schema: "https://schemas.gert.dev/runbook/v2.json"
apiVersion: runbook/v2
id: deploy-service
name: Deploy Service

x-myorg-cost-center: "engineering-ops"
x-myorg-change-id: "CHG-20260418-001"
x-otel-service-name: "payment-service"
\end{verbatim}

\paragraph{At the step level}

\begin{verbatim}
- step:
    id: deploy
    type: tool
    title: Deploy to production
    x-myorg-risk-level: high
    x-myorg-rollback-id: "{{ .deploy_id }}"
\end{verbatim}

\subsection{Extension Validation}

\begin{enumerate}
  \item \textbf{Unregistered extension fields} — The JSON~Schema validator
        accepts any field matching \texttt{\^{}x-[a-z][a-z0-9-]+\$} without
        requiring a registered schema.  This enables zero-friction annotation.
        The \texttt{gert validate} command emits a \texttt{WARNING} for each
        unregistered extension field.
  \item \textbf{Registered extension fields} — An extension may ship a JSON
        Schema fragment that constrains its field values.  Once registered, the
        validator applies the fragment and reports violations as errors, not
        warnings.
  \item \textbf{Unknown non-extension fields} — Any field that does not match
        a core field name and does not start with \texttt{x-} is a hard
        validation error (\texttt{additionalProperties: false}).
\end{enumerate}

\subsection{Extension Registration Format}

\begin{verbatim}
# .gert/extensions/myorg.yaml
apiVersion: extension/v1
namespace: myorg
description: MyOrg internal runbook annotations

fields:
  cost-center:
    type: string
    description: Cost center code for billing
    pattern: "^[a-z]+-[a-z]+"
  risk-level:
    type: string
    enum: [low, medium, high, critical]
  change-id:
    type: string
    pattern: "^CHG-[0-9]{8}-[0-9]{3}$"
  rollback-id:
    type: string

applies-to: [runbook, step]
\end{verbatim}

The extension schema is merged into the live JSON~Schema when the extension is
registered, enabling IDE auto-completion for extension fields as well.

% ============================================================
\section{Structural vs.\ Semantic Validation}
\label{sec:validation}
% ============================================================

v2 validation runs in two sequential phases.  Both must pass for a runbook to
be considered valid.

\subsection{Phase 1: Structural Validation}
\label{subsec:structural}

Structural validation is pure JSON~Schema~Draft~2020-12 evaluation.  It checks:

\begin{itemize}
  \item \textbf{Required fields} present (\texttt{id}, \texttt{name},
        \texttt{apiVersion}, \texttt{flow}).
  \item \textbf{Field types} correct (string vs.\ integer vs.\ object vs.\ array).
  \item \textbf{Enum values} valid (e.g.\ \texttt{kind} must be one of the
        four defined values).
  \item \textbf{Pattern constraints} satisfied (e.g.\ \texttt{id} must match
        slug pattern; \texttt{timeout} must match \texttt{[0-9]+(s|m|h)}).
  \item \textbf{minItems} / \textbf{minLength} constraints met
        (e.g.\ \texttt{flow} must have at least one node;
        \texttt{choices.options} must have at least two options).
  \item \textbf{additionalProperties} — no unknown core fields present.
\end{itemize}

Structural validation is performed by the generated JSON~Schema artifact
(see \S\ref{sec:schema-artifacts}).  It is stateless: it does not load tool
definitions, resolve imports, or evaluate templates.

\paragraph{Implementation note for Brian}

The Go runtime uses the \texttt{invopop/jsonschema} library to generate the
JSON~Schema from struct tags at startup (no static file required), then validates
the parsed YAML document against it.  The \texttt{yaml.v3 KnownFields(true)}
call provides an additional pre-JSON-Schema pass that rejects unknown fields
before the struct is even unmarshalled.

\subsection{Phase 2: Semantic Validation}
\label{subsec:semantic}

Semantic validation enforces domain rules that cannot be expressed in
JSON~Schema.  It runs after structural validation and requires the full runbook
AST and loaded tool/provider definitions.  Semantic checks include:

\begin{enumerate}
  \item \textbf{Step ID uniqueness} — All step IDs in the flat walk of the tree
        must be unique within the runbook.
  \item \textbf{Invoke target resolution} — Every \texttt{invoke.runbook} value
        must resolve to a declared \texttt{imports} alias or a relative file path
        that exists on disk.
  \item \textbf{Branch condition variable references} — Each \texttt{condition}
        template references only variables that are in scope at that point in the
        execution tree (declared as \texttt{inputs}, \texttt{vars}, or captured
        by an earlier step on a reachable path).
  \item \textbf{Tool reference resolution} — Every \texttt{tool.name} in a tool
        step must match a declared \texttt{toolRefs} entry, and the named
        \texttt{.tool.yaml} file must be loadable and structurally valid.
  \item \textbf{Tool action existence} — The \texttt{tool.action} must be
        declared in the referenced tool definition.
  \item \textbf{Provider binding resolution} — Every \texttt{inputs.*.from:
        <provider>.<field>} must resolve to a loaded provider whose field
        manifest declares that field name.
  \item \textbf{Circular invoke detection} — The invoke graph must be a DAG;
        any cycle is a semantic error.
  \item \textbf{Capture write-before-read} — A step that reads a captured
        variable via a template must have a reachable predecessor step that
        writes it; a warning is emitted if the write is only on some branches.
  \item \textbf{Governance policy pre-check} — The static governance rules are
        evaluated against the step contracts declared in the runbook.  Any
        step with an effect that would trigger \texttt{deny} is flagged as a
        semantic error.  Steps that would trigger \texttt{require-approval} are
        flagged as warnings.
  \item \textbf{Output declaration completeness} — Every declared \texttt{output}
        must reference a capture name that is written by at least one step.
  \item \textbf{Compensate ordering} — Each \texttt{compensate} step must appear
        after the step it is compensating for in tree walk order.
\end{enumerate}

\paragraph{Error reporting}

Both validation phases produce structured error reports with:
\begin{itemize}
  \item Error code (e.g.\ \texttt{SCH-001}, \texttt{SEM-007})
  \item File path and YAML line/column (where available)
  \item Human-readable message
  \item Suggested fix (where computable)
\end{itemize}

\begin{verbatim}
$ gert validate service-health.runbook.yaml
[SEM-003] step "check_http": tool action "head" not declared in
          tools/curl.tool.yaml (declared actions: get, post, download)
          --> service-health.runbook.yaml:34:7
[SEM-007] step "eval": template references variable "dns_result" which
          is only written on branch "DNS check" — may be unbound
          --> service-health.runbook.yaml:51:12
          hint: add a default value in vars: or guard with {{ if .dns_result }}
\end{verbatim}

% ============================================================
\section{Schema Artifacts}
\label{sec:schema-artifacts}
% ============================================================

\subsection{Shipped Artifacts}

gert v2 ships the following schema artifacts:

\begin{center}
\begin{tabular}{lp{9cm}}
\hline
\textbf{Artifact} & \textbf{Description} \\
\hline
\texttt{runbook.v2.schema.json}   & JSON~Schema Draft~2020-12 for runbook documents \\
\texttt{tool.v2.schema.json}      & JSON~Schema for tool definitions \\
\texttt{provider.v2.schema.json}  & JSON~Schema for provider definitions \\
\hline
\end{tabular}
\end{center}

All three schemas are generated at runtime from Go struct tags via
\texttt{invopop/jsonschema} and served by \texttt{gert schema export}.
They are also embedded in the binary and served at
\texttt{https://schemas.gert.dev/*} for IDE tooling.

\subsection{\texttt{gert schema export}}

\begin{verbatim}
# Export runbook schema (default)
gert schema export

# Export specific schema
gert schema export --kind runbook   # runbook | tool | provider | bundle

# Export full bundle (all three)
gert schema export --kind bundle > gert-v2-schemas.json

# Validate a document against the exported schema
gert schema export | check-jsonschema --schemafile - runbook.yaml
\end{verbatim}

The \texttt{bundle} format is a JSON object with three keys:
\texttt{runbook}, \texttt{tool}, and \texttt{provider}.

\subsection{VS Code Integration}

The gert VS~Code extension configures schema-backed YAML validation via the
\texttt{yaml-language-server} extension:

\begin{verbatim}
# .vscode/settings.json (auto-generated by gert vscode init)
{
  "yaml.schemas": {
    "https://schemas.gert.dev/runbook/v2.json": "**/*.runbook.yaml",
    "https://schemas.gert.dev/tool/v2.json":    "**/*.tool.yaml",
    "https://schemas.gert.dev/provider/v2.json": "**/*.provider.yaml"
  }
}
\end{verbatim}

Because the schema URL matches the \texttt{\$schema} field embedded in each
document, IDEs with built-in JSON~Schema support (VS~Code 1.80+, IntelliJ 2024+)
pick up validation and auto-completion without any per-project configuration.

The gert VS~Code extension additionally provides:
\begin{itemize}
  \item Step-type-aware field completion via the \texttt{schema/stepFields}
        JSON-RPC method (returns the field set for a given \texttt{type} value).
  \item Tool argument completion via \texttt{schema/toolArgs} (returns the
        \texttt{args} schema for a given tool+action pair from the live tool
        catalog).
  \item Inline governance warnings derived from \texttt{exec/dryRun} pre-check.
\end{itemize}

\subsection{Schema Version Manifest}

Each schema artifact carries a version manifest in its \texttt{\$id} and a
\texttt{x-gert-version} annotation:

\begin{verbatim}
{
  "$schema": "https://json-schema.org/draft/2020-12/schema",
  "$id": "https://schemas.gert.dev/runbook/v2.json",
  "x-gert-version": "2.0.0",
  "x-gert-released": "2026-04-18",
  "title": "gert runbook v2",
  ...
}
\end{verbatim}

Automated compatibility testing (part of the gert CI pipeline) runs every
v1 example runbook through the v2 compatibility parser and asserts no regression,
and runs every v2 example through the exported JSON~Schema validator.
